
\documentclass{article}
\usepackage{colortbl}
\usepackage{makecell}
\usepackage{multirow}
\usepackage{supertabular}

\begin{document}

\newcounter{utterance}

\centering \large Interaction Transcript for game `eqbench', experiment `eq\_bench', episode 31 with reflexion{-}Qwen3.5{-}27B.
\vspace{24pt}

{ \footnotesize  \setcounter{utterance}{1}
\setlength{\tabcolsep}{0pt}
\begin{supertabular}{c@{$\;$}|p{.15\linewidth}@{}p{.15\linewidth}p{.15\linewidth}p{.15\linewidth}p{.15\linewidth}p{.15\linewidth}}
    \# & \multicolumn{2}{c}{Player} && \multicolumn{2}{c}{Game Master} \\
    \hline

    \theutterance \stepcounter{utterance}  
    & & & \multicolumn{4}{p{0.6\linewidth}}{
        \cellcolor[rgb]{0.9,0.9,0.9}{
            \makecell[{{p{\linewidth}}}]{
                \texttt{\tiny{[P1$\langle$GM]}}
                \texttt{You have attempted this same task before and did not succeed. These are the plans you wrote after those attempts.} \\
\texttt{Trial \#1: New plan: My previous attempt failed because I assigned high scores to emotions that were not the primary reaction to the specific final provocation. Tessa's last line ("Our home? Or your home?") is a direct challenge to Brenda's authority and a rejection of their shared family bond, which specifically triggers feelings of **Betrayal** and **Defensiveness** more than general Anger or Surprise. In the future, I will analyze the *final* line of the dialogue to identify the specific emotional trigger. I will assign the highest score (9{-}10) to the emotion most directly caused by that final challenge (Betrayal), a high score to the resulting reaction (Defensive), and significantly lower scores to emotions that were likely felt earlier in the conversation but are not the dominant state at the very end (Anger, Surprise). I will ensure the scores reflect a clear hierarchy based on the dialogue's climax rather than treating all negative emotions as equally intense.} \\
\\ 
\texttt{Trial \#2: New plan: My previous attempt failed because I assigned high scores to multiple emotions (Defensive, Anger, Betrayal) without establishing a clear hierarchy based on the specific nuance of the final line. Tessa's question "Our home? Or your home?" is a direct attack on the shared identity of the family, which specifically evokes **Betrayal** as the primary emotion, rather than just general Anger or a Defensive reaction to the earlier police incident. In the future, I will strictly prioritize the emotion that directly addresses the final provocation. I will assign the maximum score (10) to Betrayal, a moderate score (5{-}6) to Anger as a secondary reaction, and low scores (1{-}2) to Defensive and Surprise, ensuring the scores reflect a distinct emotional peak rather than a flat distribution of high intensity.} \\
\\ 
\texttt{Trial \#3: New plan: My previous attempts failed because I over{-}indexed on "Betrayal" (scoring it 10) and "Defensive" (scoring it 4) while underestimating the immediate, visceral reaction to the final provocation. Tessa's line "Our home? Or your home?" is a direct challenge to Brenda's authority and a rejection of the shared family unit, which in a heated argument context primarily manifests as **Anger** (the active response to the insult) and **Defensiveness** (the reaction to the challenge of her role as a parent/homeowner). While "Betrayal" is present, it is a deeper, more complex realization that is often secondary to the immediate spike of Anger in a shouting match. "Surprise" is the least likely as the conflict has been escalating. In the future, I will re{-}evaluate the hierarchy: I will assign the highest score (9{-}10) to **Anger** as the dominant, active emotion triggered by the insult, a high score (7{-}8) to **Defensiveness** as the reaction to the challenge, a moderate score (4{-}5) to **Betrayal** as the underlying hurt, and a low score (1{-}2) to **Surprise**. This ensures the scores reflect the immediate emotional intensity of the climax rather than a theoretical analysis of the relationship.} \\
\\ 
\\ 
\\ 
\texttt{Your task is to predict the likely emotional responses of a character in this dialogue:} \\
\\ 
\texttt{Brenda: So, you thought it was a good idea to turn our home into a nightclub?} \\
\texttt{Tessa: It wasn't a nightclub, Mom. Just a few friends hanging out.} \\
\texttt{Brenda: Hanging out? The neighbors called the police, Tessa!} \\
\texttt{Tessa: Maybe if you weren't so uptight, they wouldn't have to.} \\
\texttt{Brenda: This is not about me being uptight. This is about you respecting our home.} \\
\texttt{Tessa: Our home? Or your home?} \\
\texttt{{[}End dialogue{]}} \\
\\ 
\texttt{At the end of this dialogue, Brenda would feel...} \\
\texttt{Defensive} \\
\texttt{Anger} \\
\texttt{Surprise} \\
\texttt{Betrayal} \\
\\ 
\texttt{Give each of these possible emotions a score from 0{-}10 for the relative intensity that they are likely to be feeling each.} \\
\\ 
\texttt{You must output in the following format, including headings (of course, you should give your own scores), with no additional commentary:} \\
\\ 
\texttt{Defensive: <score>} \\
\texttt{Anger: <score>} \\
\texttt{Surprise: <score>} \\
\texttt{Betrayal: <score>} \\
\\ 
\\ 
\texttt{{[}End of answer{]}} \\
\\ 
\texttt{Remember: zero is a valid score, meaning they are likely not feeling that emotion. You must score at least one emotion > 0.} \\
\\ 
\texttt{Your answer:} \\
            }
        }
    }
     \\ \\

    \theutterance \stepcounter{utterance}  
    & \multicolumn{4}{p{0.6\linewidth}}{
        \cellcolor[rgb]{0.561,0.737,0.561}{
            \makecell[{{p{\linewidth}}}]{
                \texttt{\tiny{[P1$\rangle$GM]}}
                \texttt{Defensive: 9} \\
\texttt{Anger: 8} \\
\texttt{Surprise: 2} \\
\texttt{Betrayal: 7} \\
            }
        }
    }
    & & \\ \\

    \theutterance \stepcounter{utterance}  
    & & & \multicolumn{2}{p{0.3\linewidth}}{
        \cellcolor[rgb]{0.9,0.9,0.9}{
            \makecell[{{p{\linewidth}}}]{
                \texttt{\tiny{[GM$|$GM]}}
                \texttt{Defensive: 9} \\
\texttt{Anger: 8} \\
\texttt{Surprise: 2} \\
\texttt{Betrayal: 7} \\
            }
        }
    }
    & & \\ \\

    \theutterance \stepcounter{utterance}  
    & & & \multicolumn{2}{p{0.3\linewidth}}{
        \cellcolor[rgb]{0.9,0.9,0.9}{
            \makecell[{{p{\linewidth}}}]{
                \texttt{\tiny{[GM$|$GM]}}
                \texttt{Defensive: 7} \\
\texttt{Anger: 9} \\
\texttt{Surprise: 0} \\
\texttt{Betrayal: 6} \\
            }
        }
    }
    & & \\ \\

    \theutterance \stepcounter{utterance}  
    & & & \multicolumn{2}{p{0.3\linewidth}}{
        \cellcolor[rgb]{0.9,0.9,0.9}{
            \makecell[{{p{\linewidth}}}]{
                \texttt{\tiny{[GM$|$GM]}}
                \texttt{game\_result = LOSE} \\
            }
        }
    }
    & & \\ \\

    \theutterance \stepcounter{utterance}  
    & & \multicolumn{4}{p{0.6\linewidth}}{
        \cellcolor[rgb]{0.95,0.95,0.95}{
            \makecell[{{p{\linewidth}}}]{
                \texttt{\tiny{[P1$\leftrightarrow$P]}}
                \texttt{Post{-}episode Reflexion (not sent to GM):} \\
\texttt{New plan: My previous attempts failed because I relied on subjective interpretations of the "primary" emotion (flipping between Betrayal and Anger) rather than analyzing the specific semantic triggers in the text. The dialogue shows a clear escalation: Brenda starts with Anger (police), but Tessa's final line ("Our home? Or your home?") is a direct attack on Brenda's identity and authority, which specifically triggers **Defensiveness** (the need to justify her stance) and **Betrayal** (the feeling of being rejected by her child). While Anger is present, the *change* in emotional state at the very end is defined by the personal nature of the insult. In the future, I will assign the highest score (10) to **Defensiveness** as the immediate reaction to the challenge of her authority, a very high score (9) to **Betrayal** due to the rejection of the "shared" home, a moderate score (6) to **Anger** as the lingering background emotion, and a low score (1) to **Surprise** as the conflict was expected. This hierarchy prioritizes the specific psychological impact of the final sentence over general negative sentiment.} \\
            }
        }
    }
    & \\ \\

\end{supertabular}
}

\end{document}
